\markdownRendererDocumentBegin
\markdownRendererWarning{The `hybrid` option has been soft-deprecated.}{The `hybrid` option has been soft-deprecated.}{Consider using one of the following better options for mixing TeX and markdown: `contentBlocks`, `rawAttribute`, `texComments`, `texMathDollars`, `texMathSingleBackslash`, and `texMathDoubleBackslash`. For more information, see the user manual at <https://witiko.github.io/markdown/>.}{Consider using one of the following better options for mixing TeX and markdown: `contentBlocks`, `rawAttribute`, `texComments`, `texMathDollars`, `texMathSingleBackslash`, and `texMathDoubleBackslash`. For more information, see the user manual at <https://witiko.github.io/markdown/>.}\markdownRendererSectionBegin
\markdownRendererHeadingOne{Identificar dos arboles isomorficos}\markdownRendererInterblockSeparator
{}Determinar si dos árboles son isomorfos implica verificar si existe una biyección que preserve adyacencias. El algoritmo $AHU (Aho, Hopcroft, Ullman)$ resuelve esto en $O(n)$ mediante hashing canónico: asigna etiquetas únicas a subárboles desde las hojas hacia la raíz, creando una representación invariante bajo isomorfismo. Dos árboles son isomorfos si sus representaciones canónicas coinciden, siendo útil en detección de patrones estructurales.\markdownRendererInterblockSeparator
{}\markdownRendererInputFencedCode{./_markdown_main/2f816965fbd6325825132dedc543b12e.verbatim}{cpp}{cpp}
\markdownRendererSectionEnd \markdownRendererDocumentEnd